\documentclass[10pt,a4paper]{amsart}
\usepackage[margin=19mm]{geometry}
\usepackage{amsmath,amssymb,mathtools}
\usepackage[scr=rsfs,cal=euler]{mathalfa}
\usepackage{xfrac}
\usepackage[unicode,hidelinks,bookmarks=false]{hyperref}
\newcommand{\ie}{i.e.\ }
\newcommand{\cf}{cf.\ }
\newcommand{\resp}{respectively\ }
\newcommand{\Subsection}{\subsection}
\providecommand{\nobreakdash}{\nobreak\hbox{-}}
\setlength{\emergencystretch}{2em}
\multlinegap0pt
\usepackage{amssymb}
\usepackage{float}
\usepackage{xcolor}
\usepackage{tikz}
\newtheorem{cor}{Corollary}
\newtheorem{prop}{Proposition}
\title{Rate-induced tipping in a coral reef ecosystem: A slow increase in fishing effort can induce reef collapse}
\author{Irakli Antidze, Brian Hennessy, Nikola Popovi\'c, Zak Sattar}
\date{Source: arXiv:2605.23487v1 (CC BY 4.0)}
\begin{document}
\maketitle

\noindent\emph{Source and licence.} Rate-induced tipping in a coral reef ecosystem: A slow increase in fishing effort can induce reef collapse. Version arXiv:2605.23487v1 (CC BY 4.0), distributed by arXiv under the Creative Commons Attribution 4.0 International licence, http://creativecommons.org/licenses/by/4.0/, as recorded in the arXiv licence metadata for this exact version and shown on the abstract page https://arxiv.org/abs/2605.23487v1. Full licence notice: \texttt{licenses/arxiv-2605.23487-CC-BY-4.0.txt}.\par\smallskip

\noindent\textit{Editorial scope.} Counted unit: the article's prop ordering (source0.tex lines 626--628). The article's complete original proof of it is delivered with it (source0.tex lines 629--640, one \texttt{proof} environment of 12 lines). No mathematics is written, repaired or re-derived: the statement and the proof are the article's own lines, in the article's own notation, and no retained line is substituted or re-flowed.  Retained uncounted as the source-local context the proof needs: lines 336--341; lines 363--366; lines 597--600.  Nothing was omitted from any retained span, so the package carries no omission record.  No retained line contains a citation, so the wrapper carries no bibliography and no bibliography entry.  No retained line contains a figure environment, an image-inclusion command or drawing code, so no image is delivered and the package declares no asset dependency.  Editorial formatting only, in addition: an AMS \texttt{amsart} preamble that restates the article's own declarations and macros used by the retained lines, namely the environments \texttt{cor}, \texttt{prop} on separate counters, together with the standard packages those lines need.  The retained environments print in this document as Corollary 1, Proposition 1 in order of appearance, whereas the article prints the same items with the numbering of its own full document.  The complete native source of the article as distributed in the arXiv e-print for this exact version is bundled unchanged as \texttt{sources/201213/source0.tex}.
\begin{equation}
    H_I = \frac{\frac{\beta}{\lambda}-d+\sqrt{\big(\frac{\beta}{\lambda}+d\big)^2+\frac{4\beta}{\lambda}}}{2(d+1)},\quad
    A_I = \frac{\alpha}{s(H_I)},\quad\text{and}\quad
    C_I = 1-\beta - \frac{\alpha}{s(H_I)};
    \label{eq: coexistence point}
\end{equation}

\begin{cor}
    Suppose that $\lambda d^2<1$. Then, it follows that $\alpha^{+}\leq\hat\alpha$.
    \label{cor: hat bound}
\end{cor}

\begin{equation}
    n_{0,\alpha} = \big\{(H,A,C,\alpha)\in S^2_{0,\alpha}\,\big|\, \alpha  = \lambda s(H)^2\big[H+s(H)(1+H)^2\big]\big\}.
    \label{eq: fold curve}
\end{equation}

\begin{prop} There holds $H_I>\hat{H}$ ($H_I<\hat{H}$) if and only if $\alpha^+<\hat{\alpha}<\alpha^{\ast}$ ($\alpha^{\ast}<\alpha^+<\hat{\alpha}$).
\label{prop: ordering}
\end{prop}

\begin{proof}
The fold curve $n_{0,\alpha}$, which is a one-dimensional curve in the $4$-dimensional $(H,A,C,\alpha)$-space, may be parametrised by the single variable $H$. By Equation~\eqref{eq: fold curve}, we may write
 \begin{equation}
 \begin{split}
     n_{0,\alpha}=&\biggl\{(H,A(H,\alpha(H)),C(H,\alpha(H)),\alpha(H))\in S_{0,\alpha}^{2}\,\bigg|\,  A(H,\alpha(H))=\frac{\alpha(H)}{s(H)}, \\ & C(H,\alpha(H))=1-\frac{\alpha(H)}{s(H)}-\lambda Hs(H),\ \alpha(H)=\lambda s(H)^2(H+s(H)(1+H)^2\biggr\},
     \label{eq: fold curve param}
\end{split}
 \end{equation}
with $\alpha(H)$ as given in Equation~\eqref{eq: fold curve}. In particular, since $\alpha'(H)>0$, the function $\alpha(H)$ is monotone. Also, recall that $H_I$ is independent of $\alpha$; see Equation~\eqref{eq: coexistence point}. Hence, the question of whether $H_I>\hat{H}$ or $H_I<\hat{H}$ is equivalent to considering whether the coexistence state $e_{I,\alpha}$ has crossed $n_{0,\alpha}$ before the saddle-node bifurcation occurs at $\alpha=\hat{\alpha}$. Specifically, $H_I>\hat{H}$ if and only if $e_{I,\alpha}$ has not crossed $n_{0,\alpha}$ when $\alpha=\hat{\alpha}$, which forces $\hat\alpha<\alpha^\ast$. Furthermore, by Corollary~\ref{cor: hat bound}, $\alpha^+<\hat\alpha$, and we may conclude that $H_I>\hat{H}$ if and only if $\alpha^+<\hat{\alpha}<\alpha^{\ast}$.

Similarly, if $H_I<\hat{H}$, then the coexistence state $e_{I,\alpha}$ must have crossed $n_{0,\alpha}$ before $\hat\alpha$, which is equivalent to $\alpha^{\ast}<\hat\alpha$. Furthermore, the case where $\alpha^+<\alpha^{\ast}<\hat\alpha$ is excluded, as it would then follow that $e_{I,\alpha}$ crosses $n_{0,\alpha}$ in a point with $C<0$; however, the fold curve $n_{0,\alpha}$ only lies in $\{C<0\}$ when $\alpha>\hat\alpha$. In particular, $H_I<\hat{H}$ if and only if $\alpha^{\ast}<\alpha^+<\hat\alpha$. 
\end{proof}

\end{document}
